\documentclass[12pt,a4paper]{article}

% Paquetes para español
\usepackage[utf8]{inputenc}
\usepackage[T1]{fontenc}
% Babel with Spanish support would go here if available

% Paquetes adicionales
\usepackage{geometry}
\usepackage{setspace}
\usepackage{csquotes}
\usepackage{hyperref}

% Configuración de márgenes
\geometry{
    left=2.5cm,
    right=2.5cm,
    top=2.5cm,
    bottom=2.5cm
}

% Espaciado
\setstretch{1.15}

% Información del documento
\title{\textbf{EL CATÁLOGO DE LA BIBLIOTECA DE FERNÁN DE VELASCO:} \\ 
\textbf{UN TESORO BIBLIOGRÁFICO DEL SIGLO XVIII} \\
\large Estudio filológico, bibliotecónomico y bibliográfico de un catálogo excepcional}
\author{}
\date{}

\begin{document}

\maketitle

\begin{abstract}
El presente trabajo analiza el catálogo de la biblioteca de D. Fernán de Velasco, miembro del Consejo y Cámara de Su Majestad, tasado en Madrid en 1791 por Antonio Baylo. Esta colección bibliográfica, actualmente custodiada en la Biblioteca Nacional de España, constituye un testimonio excepcional de la cultura libresca española del siglo XVIII y representa una fuente primaria invaluable para comprender los hábitos de lectura, las redes comerciales del libro y las corrientes intelectuales de la Ilustración española.
\end{abstract}

\section{INTRODUCCIÓN}

El catálogo de la Casa de Velasco es un documento manuscrito en dos tomos en folio que registra meticulosamente una de las bibliotecas privadas más notables del Madrid dieciochesco. Tasado el 4 de octubre de 1791, el conjunto alcanzó la considerable suma de 232,951 reales y 8 maravedís, lo que evidencia tanto la amplitud de la colección como su valor económico y cultural.

El análisis filológico y bibliotecónomico de este catálogo revela no solo el contenido de una biblioteca aristocrática, sino también los criterios de organización, valoración y clasificación bibliográfica vigentes en el último tercio del siglo XVIII español.

\section{METODOLOGÍA CATALOGRÁFICA}

\subsection{Estructura de las entradas}

Cada entrada del catálogo sigue un sistema normalizado que incluye:

\begin{enumerate}
\item[a)] \textbf{Autor:} Generalmente en latín o en la lengua original de la obra, con variantes como ``Eiusd.'' (\textit{Eiusdem} = del mismo autor) o ``Id.'' (\textit{Idem} = lo mismo) cuando se registran varias obras del mismo autor consecutivamente.

\item[b)] \textbf{Título:} Puede aparecer en latín, castellano, italiano, francés o portugués, respetando la lengua de publicación.

\item[c)] \textbf{Formato:} Expresado en sistema de plegado de hojas:
\begin{itemize}
    \item fol. (folio): el formato mayor
    \item 4º (cuarto): hoja plegada dos veces
    \item 8º (octavo): hoja plegada tres veces
    \item 12º (dozavo): hoja plegada cuatro veces
    \item 16º (dieciseisavo): el formato menor
\end{itemize}

\item[d)] \textbf{Lugar y año de impresión:} Con abreviaturas geográficas típicas de la época (Mad. = Madrid, Lugd. = Lugduni/Lyon, Antuerp. = Amberes, etc.)

\item[e)] \textbf{Número interno:} Sistema de catalogación particular de la biblioteca, precedido por ``N.'' o ``Numero''

\item[f)] \textbf{Precio:} Expresado en reales y maravedís, con notaciones como ``r'' o ``s'' (reales), indicando el valor de tasación de 1791.
\end{enumerate}

\subsection{Peculiaridades del sistema de registro}

Un rasgo distintivo del catálogo es la omisión del autor en entradas subsecuentes cuando varias obras del mismo pertenecen a un único autor. Esta práctica, común en la catalogación manuscrita de la época, requiere una lectura contextual para identificar correctamente la autoría.

Ejemplos representativos:

\begin{quote}
``Carthagena (P. Joannis) De Jure belli Romani Pontificis adversum Ecclesiae Inviolantes 8º Romae 1609 N. 190 6 \\
Carthier (P. Galli) Immortalitas animae nostrae adversus nostrae aetatis Philosophos 8º Augustae Vindelicor. 1767 N. 83 6''
\end{quote}

En este caso, cada entrada tiene su autor explícito. Sin embargo, cuando aparece:

\begin{quote}
``Velasco (Alvaro) De Jure Emphyteutico fol. Olisiponae 1614. Num. 246 15r \\
Velasco Souvea (F. Francis) Aclamacion Do Rey de Portugal Dom Joao 1º fol. Lisboa 1644 306 30r \\
Id. Lexidia de Alemania y D. Castilla en la Prisión del Infante de Portugal D. Duarte fol. Lisboa 1652 306 32r''
\end{quote}

El ``Id.'' indica que la tercera obra también pertenece a Velasco Souvea.

\section{ANÁLISIS FILOLÓGICO: DIVERSIDAD LINGÜÍSTICA Y CULTURAL}

\subsection{Composición políglota de la colección}

El catálogo revela una biblioteca profundamente cosmopolita y erudita. Los idiomas representados incluyen:

\begin{itemize}
    \item \textbf{Latín:} La lengua predominante, especialmente en obras jurídicas, teológicas y filosóficas
    \item \textbf{Castellano:} Ampliamente presente en obras históricas, literarias y administrativas
    \item \textbf{Portugués:} Notable presencia, reflejo de las estrechas relaciones ibéricas
    \item \textbf{Italiano:} Obras literarias y políticas
    \item \textbf{Francés:} Textos ilustrados y obras contemporáneas
    \item \textbf{Griego y Hebreo:} Textos religiosos y gramaticales especializados
    \item \textbf{Árabe:} Gramáticas y vocabularios, testimonio del interés orientalista
\end{itemize}

Esta diversidad lingüística evidencia el perfil de un lector culto, vinculado al aparato estatal y con intereses intelectuales amplios, característico de la élite ilustrada española.

\subsection{Análisis de procedencias geográficas}

El estudio de los lugares de impresión revela las principales rutas del comercio libresco europeo que abastecían el mercado español:

\subsubsection*{Centros de impresión españoles:}
\begin{itemize}
    \item \textbf{Madrid} (Mad., Matriti): Principal centro editorial
    \item \textbf{Sevilla:} Importante para obras americanistas
    \item \textbf{Barcelona, Valencia, Salamanca, Alcalá:} Centros universitarios y eclesiásticos
\end{itemize}

\subsubsection*{Centros europeos:}
\begin{itemize}
    \item \textbf{Lyon} (Lugduni, Lugd.): Principal proveedor francés de libros latinos
    \item \textbf{Amberes} (Antuerp., Antuerpiae): Crucial en el comercio con los Países Bajos
    \item \textbf{Venecia} (Venetiis): Obras italianas y clásicas
    \item \textbf{París} (Parisiis): Textos ilustrados y contemporáneos
    \item \textbf{Lisboa} (Olisiponae): Literatura lusa e historia colonial
    \item \textbf{Roma} (Romae): Textos teológicos y oficiales
\end{itemize}

Esta geografía editorial revela las redes comerciales y culturales que conectaban España con el resto de Europa durante el siglo XVIII.

\section{ANÁLISIS BIBLIOTECÓNOMICO: CAMPOS TEMÁTICOS}

\subsection{Derecho y jurisprudencia}

Una parte sustancial de la biblioteca está dedicada al derecho, tanto civil como canónico. Obras fundamentales incluyen:

\begin{itemize}
    \item Comentarios al Derecho Romano (múltiples ediciones de Vinnius)
    \item Derecho Canónico (Carthagena, Casaubonius)
    \item Derecho Real Español (Nueva Recopilación de las Leyes de España)
    \item Tratados jurídicos especializados sobre temas mercantiles, marítimos y coloniales
\end{itemize}

Esta abundancia jurídica concuerda con la posición de Fernán de Velasco en el Consejo y Cámara Real.

\subsection{Historia y genealogía}

El catálogo refleja un interés profundo por:

\begin{itemize}
    \item Historia de España (Mariana, Cascales, Colmenares)
    \item Historia local y regional (Toledo, Murcia, Segovia, Valencia)
    \item Genealogías nobiliarias (Casa de Velasco, casas señoriales)
    \item Historia portuguesa y colonial
    \item Historia eclesiástica
\end{itemize}

Ejemplos notables:

\begin{quote}
``Fernandez de Velasco (D. Bernardino) Advertencias y ala Historia de España del Padre Mariana 4º ciud 1613 Nº 134 12''
\end{quote}

\begin{quote}
``Velasco (D. Pedro Fernandez) Segundo Libro de Sevilla, con una Sumaria Relacion del Linaje de Velasco por Pedro Cintiano f. Milan 1671 V. 198 20r''
\end{quote}

\subsection{Teología y religión}

Abundan obras de:
\begin{itemize}
    \item Teología escolástica
    \item Vidas de santos
    \item Constituciones sinodales
    \item Literatura devocional
    \item Controversias religiosas
\end{itemize}

\subsection{Filología y humanidades}

Notable presencia de:

\begin{itemize}
    \item Gramáticas (latina, árabe, hebrea, japonesa)
    \item Diccionarios y vocabularios multilingües
    \item Ediciones clásicas (Virgilio, especialmente traducciones al castellano)
    \item Retórica y poética
    \item Emblemática (Alciato)
\end{itemize}

Ejemplo destacado:

\begin{quote}
``Vocabulario dela Lengua Arábiga, en lengua Castellana: 4.° 1645. Num. 31 36''
\end{quote}

\begin{quote}
``Collado (J. Didaci) Grammatica Linguae Japonicae 4° Romae 1632 N. 303 8''
\end{quote}

\subsection{Literatura}

\begin{itemize}
    \item Obras poéticas españolas e italianas
    \item Novelas (Mateo Alemán, Cervantes probablemente)
    \item Teatro castellano
    \item Traducciones de clásicos
\end{itemize}

\subsection{Ciencias y artes}

\begin{itemize}
    \item Tratados militares y de fortificación
    \item Medicina
    \item Matemáticas y astronomía
    \item Arte de escribir y caligrafía
    \item Geografía y cartografía
\end{itemize}

\section{ANÁLISIS ECONÓMICO: EL MERCADO DEL LIBRO}

\subsection{Sistema de valoración}

Los precios registrados en reales permiten analizar el mercado del libro antiguo en 1791:

\subsubsection*{Obras de bajo precio (menos de 10 reales):}
\begin{itemize}
    \item Libros devocionales pequeños
    \item Opúsculos y textos breves
    \item Obras en formato menor (12º, 16º)
\end{itemize}

\subsubsection*{Obras de precio medio (10-50 reales):}
\begin{itemize}
    \item Monografías en 4º u 8º
    \item Obras históricas y jurídicas estándar
    \item Literatura clásica
\end{itemize}

\subsubsection*{Obras de alto precio (más de 50 reales):}
\begin{itemize}
    \item Grandes folios con ilustraciones
    \item Obras en múltiples tomos
    \item Ediciones lujosas con encuadernación especial (``Pasta'', ``tafilete'')
    \item Biblias y obras litúrgicas monumentales
\end{itemize}

Ejemplo de obra valiosa:

\begin{quote}
``Biblia Sacra Hebraice et Graece cum latine interlineari interpretatione Pagnini et Benedicti Arias Montano. folio Antwerp. 1584 Nº 96. 240''
\end{quote}

\subsection{El tasador: Antonio Baylo}

El documento final del catálogo incluye la certificación del tasador:

\begin{quote}
``Antonio Baylo Madrid, 31 de junio de 1791 Dijo yo el abajo firmado que he visto reconocido y tasado los libros contenidos en este catalogo 2 tomos en folio y fueron del Sr. D. Fernan de Velasco del Consejo y Camara de S.M. A lo qual se ha dado el valor segun mi respectivo estado y mi saber y entender.''
\end{quote}

Antonio Baylo actuó como experto bibliófilo y tasador oficial, aplicando criterios de valoración basados en:
\begin{itemize}
    \item Estado de conservación
    \item Rareza y demanda
    \item Formato y calidad de impresión
    \item Encuadernación
    \item Completitud de ejemplares
\end{itemize}

\section{CONTEXTO HISTÓRICO: LA BIBLIOTECA ILUSTRADA}

\subsection{Perfil del propietario}

Fernán de Velasco, miembro del Consejo y Cámara de Su Majestad, representa el tipo del funcionario-erudito característico de la Ilustración española. Su biblioteca revela:

\begin{itemize}
    \item Formación jurídica sólida
    \item Interés genealógico (vinculado a su linaje nobiliario)
    \item Curiosidad intelectual cosmopolita
    \item Compromiso con el servicio público
    \item Participación en las redes culturales ilustradas
\end{itemize}

\subsection{El destino del catálogo}

El hecho de que el catálogo se conserve en la Biblioteca Nacional de España sugiere que, tras la muerte de Velasco, la colección pudo dispersarse o integrarse en fondos públicos, práctica común en el proceso de secularización y nacionalización de bienes culturales que caracterizó el tránsito del Antiguo Régimen al liberalismo.

\section{VALOR COMO FUENTE HISTÓRICA}

Este catálogo constituye una fuente primaria excepcional para:

\begin{enumerate}
\item[a)] \textbf{Historia del libro y la lectura:} Revela qué se leía, en qué idiomas, y qué se valoraba en la España del XVIII.

\item[b)] \textbf{Historia cultural e intelectual:} Permite reconstruir las bibliotecas de la élite funcionarial ilustrada.

\item[c)] \textbf{Historia de las mentalidades:} Los temas y autores seleccionados reflejan los intereses, preocupaciones y horizontes intelectuales de un sector social específico.

\item[d)] \textbf{Historia económica:} Los precios documentan el mercado del libro antiguo y su valoración.

\item[e)] \textbf{Filología histórica:} Las variantes ortográficas, las abreviaturas y la terminología técnica son testimonio de usos lingüísticos de época.

\item[f)] \textbf{Historia de las bibliotecas:} Documenta prácticas de catalogación, organización y tasación bibliográfica pre-modernas.
\end{enumerate}

\section{CONCLUSIONES}

El Catálogo de la Biblioteca de Fernán de Velasco es mucho más que un inventario de libros. Representa una ventana privilegiada a la cultura libresca de la Ilustración española, revelando:

\begin{enumerate}
    \item La existencia de bibliotecas privadas extraordinariamente ricas y diversas en el Madrid del siglo XVIII.
    
    \item La integración de la élite española en las corrientes intelectuales europeas, evidenciada por la presencia de obras en múltiples lenguas y de diversos centros editoriales.
    
    \item La pervivencia de intereses tradicionales (teología, derecho, genealogía) junto a preocupaciones ilustradas (ciencias, literatura contemporánea, geografía).
    
    \item La profesionalización de los oficios relacionados con el libro (tasadores, catalogadores, libreros).
    
    \item El valor económico considerable del patrimonio bibliográfico privado.
\end{enumerate}

Este documento, preservado en la Biblioteca Nacional de España, merece ser estudiado en profundidad no solo como catálogo bibliográfico, sino como artefacto cultural que condensa las aspiraciones intelectuales, los horizontes de lectura y las redes de circulación del libro en la España de finales del Antiguo Régimen.

La meticulosidad de Antonio Baylo en su tasación y el cuidado en registrar autores, títulos, formatos, lugares de impresión y precios hacen de este catálogo una herramienta indispensable para historiadores del libro, filólogos, historiadores culturales y todos aquellos interesados en comprender cómo se leía, qué se valoraba y cómo circulaban los libros en la España ilustrada.

\section*{BIBLIOGRAFÍA RECOMENDADA}

\begin{itemize}
    \item Aguilar Piñal, Francisco. \textit{La biblioteca de Jovellanos (1778)}. Madrid: CSIC, 1984.
    \item Bouza, Fernando. \textit{Corre manuscrito: Una historia cultural del Siglo de Oro}. Madrid: Marcial Pons, 2001.
    \item Domergue, Lucienne. \textit{Tres calas en la censura dieciochesca}. Toulouse: Université de Toulouse-Le Mirail, 1981.
    \item García Ejarque, Luis. \textit{La Real Biblioteca de S.M. y su personal (1712-1836)}. Madrid: Fundación Germán Sánchez Ruipérez, 1997.
    \item Infantes, Víctor, François Lopez y Jean-François Botrel (dirs.). \textit{Historia de la edición y de la lectura en España, 1472-1914}. Madrid: Fundación Germán Sánchez Ruipérez, 2003.
    \item López, François. ``Gentes y oficios de la librería española a mediados del siglo XVIII''. \textit{Nueva Revista de Filología Hispánica}, 33:1 (1984), pp. 165-185.
    \item Martínez Martín, Jesús A. (dir.). \textit{Historia de la edición en España (1836-1936)}. Madrid: Marcial Pons, 2001.
    \item Pedraza Gracia, Manuel José. \textit{Lectores y lecturas en Zaragoza (1501-1521)}. Zaragoza: Prensas Universitarias de Zaragoza, 1998.
\end{itemize}

\section*{ANEXO: EJEMPLOS REPRESENTATIVOS DE ENTRADAS}

A continuación se transcriben algunas entradas representativas que ilustran la diversidad y riqueza del catálogo:

\subsection*{[Derecho Romano]}
\begin{quote}
``Vinnii (Arnoldi) Ejusd. In quatuor libros imp. instit. commentarius: 4.° 2 tom. en un volumen Amsterdam 1665. c. N° 187. 24''
\end{quote}

\subsection*{[Historia de España]}
\begin{quote}
``Fernandez de Velasco (D. Bernardino) Advertencias y ala Historia de España del Padre Mariana 4° ciud 1613 N° 134 12''
\end{quote}

\subsection*{[Literatura clásica traducida]}
\begin{quote}
``Gregorio Hernández de Velasco In Eneida traducida en verso castellano 1615 8.º cuad. Atado 10 10s''
\end{quote}

\subsection*{[Gramática oriental]}
\begin{quote}
``Collado (J. Didaci) Grammatica Linguae Japonicae 4° Romae 1632 N. 303 8''
\end{quote}

\subsection*{[Historia local]}
\begin{quote}
``Colmenares (Diego de) Historia de la Ciudad de Segovia, y de Castilla fol. Alad. 1640 - Numero 143 60''
\end{quote}

\subsection*{[Emblemática]}
\begin{quote}
``Alciati (Andreae) Emblemata. Comment. et notis variarum 4º Patavii 1621. Nº 14 36''
\end{quote}

\subsection*{[Obras ilustradas]}
\begin{quote}
``Vegetii (Flavii) De Re Militari 4.º Antuerpiae 1585. figuras. c. 68 75r''
\end{quote}

\subsection*{[Colecciones de gran valor]}
\begin{quote}
``Biblia Sacra Hebraice et Graece cum latine interlineari interpretatione Pagnini et Benedicti Arias Montano. folio Antwerp. 1584 Nº 96. 240''
\end{quote}

\vspace{1em}
Estas entradas ejemplifican la metodología catalográfica, la diversidad temática y lingüística, y el sistema de valoración que caracteriza este extraordinario documento bibliográfico del siglo XVIII español.

\end{document}
